\documentclass[11pt,a4paper]{article}

\usepackage[a4paper,margin=1in]{geometry}
\usepackage[T1]{fontenc}
\usepackage[utf8]{inputenc}
\usepackage{lmodern}
\usepackage{microtype}
\usepackage{graphicx}
\usepackage{booktabs}
\usepackage{array}
\usepackage{xcolor}
\usepackage{hyperref}
\usepackage{listings}
\usepackage{enumitem}
\usepackage{float}
\usepackage{fancyhdr}
\usepackage{titlesec}

\definecolor{codebg}{RGB}{247,247,247}
\definecolor{linkblue}{RGB}{35,90,150}

\hypersetup{
    colorlinks=true,
    linkcolor=linkblue,
    urlcolor=linkblue,
    citecolor=linkblue,
    pdftitle={Team 1 - Private Network Service Platform - Phase 1},
    pdfauthor={Team 1}
}

\lstset{
    backgroundcolor=\color{codebg},
    basicstyle=\ttfamily\small,
    frame=single,
    breaklines=true,
    columns=fullflexible,
    showstringspaces=false
}

\pagestyle{fancy}
\fancyhf{}
\lhead{Team 1}
\rhead{Computer Networks -- Phase 1}
\cfoot{\thepage}

\titleformat{\section}{\Large\bfseries}{\thesection}{0.7em}{}
\titleformat{\subsection}{\large\bfseries}{\thesubsection}{0.7em}{}

\title{
    \vspace{2cm}
    {\LARGE\bfseries Computer Networks}\\[0.4cm]
    {\Large Private Network Service Platform}\\[0.7cm]
    {\large Phase 1 -- Build and Observe}\\[1.5cm]
    \rule{\textwidth}{0.5pt}\\[0.4cm]
    \large Team 1
    \rule{\textwidth}{0.5pt}
}
\author{
    \begin{tabular}{ll}
    \textbf{Manas Selukar} & 2401010259 \\
    \textbf{Anant Singh} & 2401010067 \\
    \textbf{Syed Darain Qamar} & 2401010472 \\
    \textbf{Akhil Sharma} & 2401020084
    \end{tabular}
    \\[0.6cm]
    \textit{Team Roles:}\\
    Manas -- Private DNS \quad
    Anant -- Edge / nginx \quad
    Akhil -- Backend A \quad
    Syed Darain -- Backend B
}
\date{Phase 1 Submission}

\begin{document}

\maketitle
\thispagestyle{empty}

\begin{abstract}
This report documents the Phase 1 implementation of a four-node private
network service platform. The system uses a macOS LAN, centralized private
DNS with dnsmasq, an nginx edge node providing HTTPS termination and
round-robin load balancing, and two simple REST backends. The implementation
is verified using \texttt{dig}, \texttt{curl}, and Wireshark captures covering
DNS, TCP, TLS, and encrypted application traffic. Controlled failure
scenarios are also documented to demonstrate layer-by-layer diagnosis.
\end{abstract}

\tableofcontents
\newpage

\section{Project Overview}

The project implements a fully local service environment on four macOS
laptops connected to the same private Wi-Fi/LAN. The objective is to show
how a client request travels through DNS resolution, TCP connection
establishment, TLS negotiation, HTTP processing, reverse proxying, and
backend selection.

The architecture deliberately keeps the application layer simple. The
network infrastructure is the primary focus.

\section{Architecture}

\subsection{Topology}

The four machines have the following roles:

\begin{itemize}[leftmargin=*]
    \item Mac 1 provides private DNS using dnsmasq and also acts as a test client.
    \item Mac 2 provides the nginx edge, reverse proxy, HTTPS termination, and load balancing.
    \item Mac 3 provides Backend A on TCP port 3001.
    \item Mac 4 provides Backend B on TCP port 3002.
\end{itemize}

The final topology diagram is stored in \texttt{docs/topology.png}.

\begin{figure}[H]
    \centering
    \IfFileExists{topology.png}{
        \includegraphics[width=0.95\textwidth]{topology.png}
    }{
        \fbox{\parbox{0.9\textwidth}{\centering
        Topology diagram: place \texttt{topology.png} in the same directory as this report before compiling.}}
    }
    \caption{Team 1 private LAN topology.}
\end{figure}

\subsection{Machine and Service Inventory}

\begin{table}[H]
\centering
\begin{tabular}{p{2.2cm}p{3.2cm}p{2.7cm}p{3.2cm}}
\toprule
\textbf{Machine} & \textbf{Role} & \textbf{Private IP} & \textbf{Services} \\
\midrule
Mac 1 -- Manas & Private DNS + Test Client & 10.7.17.68 & dnsmasq :53 \\
Mac 2 -- Anant & Edge / Reverse Proxy / Load Balancer & 10.7.21.52 & nginx :80 / :443 \\
Mac 3 -- Akhil & Backend A & 10.7.4.37 & REST API :3001 \\
Mac 4 -- Darain & Backend B + Test Client & 10.7.13.20 & REST API :3002 \\
\bottomrule
\end{tabular}
\caption{Machine and service inventory.}
\end{table}

\noindent
The complete subnet, gateway, interface, and MAC-address inventory should
be maintained alongside this report once all team members have supplied
their final network details.

\subsection{Request Flow}

The request path is:

\begin{lstlisting}
Client
  |
  | DNS query: app.team1.test (UDP/53)
  v
Mac 1 -- dnsmasq (10.7.17.68)
  |
  | DNS response: 10.7.21.52
  v
Mac 2 -- nginx Edge (10.7.21.52:443)
  |
  | TLS + HTTP
  v
+----------------------+
| nginx load balancer  |
+----------+-----------+
           |
       +---+---+
       |       |
       v       v
 Backend A  Backend B
 :3001      :3002
\end{lstlisting}

The client therefore uses the private service name and the edge node as
the public entry point. Backend IP addresses are not part of the client
URL.

\section{Private LAN}

All participating machines are intended to operate on the same private
Wi-Fi/LAN. The project requires each machine's private IPv4 address,
subnet/prefix, default gateway, active interface, and MAC address to be
recorded, together with pairwise reachability tests.

The currently recorded addresses are:

\begin{itemize}
    \item Manas: \texttt{10.7.17.68/19}, gateway \texttt{10.7.0.1}, interface \texttt{en0}.
    \item Anant: \texttt{10.7.21.52/19}, gateway \texttt{10.7.0.1}.
    \item Akhil: \texttt{10.7.4.37}.
    \item Darain: \texttt{10.7.13.20}.
\end{itemize}

The final submission should include the completed four-machine inventory
and pairwise ping evidence in \texttt{evidence/network/}.

\section{Private DNS}

\subsection{dnsmasq Configuration}

Mac 1 runs dnsmasq as the team's private DNS resolver. The repository copy
is stored at:

\begin{lstlisting}
config/dnsmasq/dnsmasq.conf
\end{lstlisting}

The private records are:

\begin{table}[H]
\centering
\begin{tabular}{ll}
\toprule
\textbf{Hostname} & \textbf{Address} \\
\midrule
\texttt{app.team1.test} & 10.7.21.52 \\
\texttt{api.team1.test} & 10.7.21.52 \\
\bottomrule
\end{tabular}
\caption{Private DNS records.}
\end{table}

dnsmasq listens on localhost and Mac 1's LAN address:

\begin{lstlisting}
listen-address=127.0.0.1,10.7.17.68
bind-interfaces
\end{lstlisting}

Public DNS queries are forwarded to configured upstream resolvers.

\subsection{Client Resolver Configuration}

At least two other Macs are configured to use Mac 1
(\texttt{10.7.17.68}) as their DNS resolver. DNS resolution was verified
using \texttt{dig}.

Example:

\begin{lstlisting}
dig app.team1.test
\end{lstlisting}

The expected private answer is:

\begin{lstlisting}
10.7.21.52
\end{lstlisting}

DNS evidence is stored in \texttt{evidence/dns/}.

\section{Backend Services}

Mac 3 and Mac 4 provide simple HTTP REST services.

\begin{table}[H]
\centering
\begin{tabular}{llll}
\toprule
\textbf{Backend} & \textbf{Host} & \textbf{Port} & \textbf{Identifier} \\
\midrule
A & 10.7.4.37 & 3001 & \texttt{X-Backend: A} \\
B & 10.7.13.20 & 3002 & \texttt{X-Backend: B} \\
\bottomrule
\end{tabular}
\caption{Backend services.}
\end{table}

Both services expose:

\begin{itemize}
    \item \texttt{GET /} -- confirms that the backend is running.
    \item \texttt{GET /api/status} -- returns JSON containing a backend identifier and status.
\end{itemize}

The backends are required to listen on a LAN-accessible interface rather
than only \texttt{127.0.0.1}.

\section{nginx Edge and Load Balancing}

Mac 2 runs nginx as the single public entry point. The upstream contains:

\begin{lstlisting}
upstream backend_servers {
    server 10.7.4.37:3001;
    server 10.7.13.20:3002;
}
\end{lstlisting}

The application and API hostnames are served over HTTPS. nginx terminates
TLS and proxies the request to the upstream backend pool.

Repeated requests were used to demonstrate responses from both backend
instances, with the \texttt{X-Backend} response header identifying the
selected backend.

\section{HTTPS and TLS}

The edge service terminates HTTPS on TCP port 443. A certificate covering
the private hostnames is installed and trusted on the client used for
verification.

The final verification uses normal certificate validation rather than
bypassing validation with \texttt{curl -k}.

The observed TLS connection included TLS 1.3 and successful certificate
verification.

The expected handshake sequence is:

\begin{lstlisting}
ClientHello
ServerHello
Certificate
Key Exchange
Finished
\end{lstlisting}

Wireshark evidence for the TLS handshake and encrypted application data
is stored in \texttt{evidence/wireshark/}.

\section{HTTP Caching}

The application response includes HTTP caching headers. An observed
response contains:

\begin{lstlisting}
Cache-Control: max-age=60
ETag: W/"..."
\end{lstlisting}

The header evidence is stored in \texttt{evidence/https/}.

\textbf{Submission note:} the Phase 1 specification requires a repeated
request showing either cache reuse or a conditional request returning
\texttt{304 Not Modified}. This demonstration should be completed and
added to the evidence before the Phase 1 review.

\section{Protocol Analysis}

\subsection{DNS}

The client first sends a DNS query for \texttt{app.team1.test}. The
team DNS server responds with the private edge address
\texttt{10.7.21.52}.

\subsection{TCP}

After DNS resolution, the client establishes a TCP connection to
\texttt{10.7.21.52:443}. The three-way handshake is:

\begin{lstlisting}
SYN
SYN-ACK
ACK
\end{lstlisting}

The client uses an ephemeral source port and the server uses the
well-known HTTPS destination port 443.

\subsection{TLS}

TLS is negotiated over the established TCP connection. The certificate
is validated by the client and subsequent application data is encrypted.

\subsection{HTTP}

The HTTP request is transmitted inside the established TLS connection.
nginx receives the request and selects one backend from the upstream pool.

\subsection{Protocol and Port Summary}

\begin{table}[H]
\centering
\begin{tabular}{llll}
\toprule
\textbf{Function} & \textbf{Protocol} & \textbf{Transport} & \textbf{Port} \\
\midrule
Private DNS & DNS & UDP & 53 \\
HTTP redirect & HTTP & TCP & 80 \\
Secure entry & HTTPS & TCP & 443 \\
Backend A & HTTP & TCP & 3001 \\
Backend B & HTTP & TCP & 3002 \\
\bottomrule
\end{tabular}
\caption{Service ports used by the system.}
\end{table}

\section{Wireshark Evidence}

The Phase 1 packet-analysis requirement is addressed using saved
Wireshark captures.

\begin{itemize}
    \item DNS query and response for \texttt{app.team1.test}.
    \item TCP SYN, SYN-ACK, and ACK sequence.
    \item TLS handshake packets.
    \item Encrypted TLS application data.
\end{itemize}

The corresponding evidence is stored under:

\begin{lstlisting}
evidence/wireshark/
\end{lstlisting}

The captures demonstrate that the application request is not visible as
plaintext HTTP payload after TLS encryption.

\section{Controlled Failure Demonstrations}

The Phase 1 failure requirements were tested deliberately.

\subsection{Wrong DNS Server}

The client was temporarily configured to use an external DNS server.
The private hostname failed to resolve, demonstrating that IP
connectivity and DNS resolution are separate layers.

\subsection{Incorrect DNS Record}

The private DNS record was temporarily changed to an incorrect
destination IP. DNS resolution still succeeded, but the client was sent
to the wrong destination and the application connection failed.

\subsection{Backend A Unavailable}

Backend A was stopped while Backend B remained available. The edge
continued to receive traffic and responses from Backend B were observed.
The final submission should demonstrate clean failover behavior without
an avoidable long timeout.

\subsection{Both Backends Unavailable}

Both backend services were stopped while the edge remained reachable.
The DNS and TLS layers could still be reached, while the application
layer could not obtain a backend response.

The Phase 1 specification expects a clear upstream failure such as
\texttt{502 Bad Gateway}; the final evidence should show the resulting
nginx behavior explicitly.

\subsection{Wrong Destination Port}

A connection was attempted against an incorrect HTTPS destination port.
DNS still resolved the hostname to the edge IP, but the TCP connection
to the incorrect service port failed.

Failure evidence is stored in:

\begin{lstlisting}
evidence/failures/
\end{lstlisting}

\section{Evidence Map}

The repository is organized so that evidence can be located quickly.

\begin{table}[H]
\centering
\begin{tabular}{p{3.2cm}p{8cm}}
\toprule
\textbf{Directory} & \textbf{Contents} \\
\midrule
\texttt{evidence/network/} & IP configuration and LAN reachability \\
\texttt{evidence/dns/} & dnsmasq status and DNS resolution \\
\texttt{evidence/https/} & HTTPS, certificate, headers, caching, load balancing \\
\texttt{evidence/wireshark/} & DNS, TCP, TLS, and encrypted traffic captures \\
\texttt{evidence/failures/} & Controlled failure demonstrations \\
\bottomrule
\end{tabular}
\caption{Phase 1 evidence organization.}
\end{table}

\section{Phase 1 Status}

The core Phase 1 infrastructure is operational:

\begin{itemize}
    \item Private LAN addressing and service roles established.
    \item Centralized private DNS configured with dnsmasq.
    \item Private application hostnames resolve to the nginx edge.
    \item nginx provides the single HTTPS entry point.
    \item Two backend services operate on fixed LAN ports.
    \item Round-robin backend selection has been demonstrated.
    \item TLS certificate validation succeeds without disabling verification.
    \item DNS, TCP, TLS, and encrypted application traffic have been captured.
    \item Required DNS and network failure scenarios have been tested.
\end{itemize}

Before the Phase 1 review, the remaining evidence items are to complete
the HTTP cache-hit or \texttt{304 Not Modified} demonstration, finalize
the four-machine network inventory and pairwise ping evidence, and ensure
the backend-failure demonstrations produce clear, evaluator-visible
results.

\section{Conclusion}

Phase 1 establishes a complete local request path from private DNS
resolution through TCP, TLS, HTTPS, nginx load balancing, and backend
response. The system is intentionally structured so that each protocol
layer can be observed independently using standard networking tools.

The resulting architecture provides a practical demonstration of how
DNS service discovery, secure ingress, reverse proxying, load balancing,
HTTP caching, and packet-level analysis work together in a private
network environment.

\end{document}
